\documentclass[11pt]{article}
\usepackage[a4paper,margin=26mm]{geometry}
\usepackage{amsmath,amssymb}
\usepackage[hidelinks]{hyperref}
\setlength{\emergencystretch}{3em}
\title{Conjecture 00000001785: the printed formula is not a step function}
\author{lizaixi01, with OpenAI Codex assistance}
\date{23 October 2026}
\begin{document}
\maketitle
\section{Literal necessary conditions}
Both versions of the source call the proposed optimal constant a
\emph{step function} of $c$, and give the explicit formula
\[
 C(c)=F(c):=\max(1,2c)\,\pi.
\]
These two necessary clauses are incompatible on the positive real
parameter window $I=(1/2,1)$. On that window the formula is
$F(c)=2\pi c$, which is injective and has uncountably many values.
A step function represented by finitely many interval indicators has
only finitely many values. Even a countable interval partition into
constant pieces has only countably many values.

We formalize both actual interval representations and their image
bounds, rather than defining ``step'' to mean countable range.
The standard finite definition is stated explicitly in Ryzhik's
Stanford analysis course~\cite[Problem 1(i), p.~1]{Ryzhik}: a finite
linear combination of indicators of disjoint intervals.
The separate countable-piece result covers a more permissive
piecewise-constant reading as well.

This is a disproof of the source's literal necessary conjunction.
It does not calculate any actual Ohsawa--Takegoshi optimal constant,
assign that constant a value, or refute an independent extension
theorem. No convention about the extension's norm or normalization
is needed: \emph{no} real function $C$ can satisfy these two clauses
on $I$. Replacing the word ``step'' by ``piecewise linear'' would
give a different assertion, which this proof does not claim to
refute. The word ``transition'' is not assumed to mean a jump; a
change of slope alone would not be a contradiction.

\section{A genuine finite interval representation}
For a set $A\subseteq\mathbb R$, the predicate
\texttt{FiniteIntervalStepOn A C} requires actual interval sets
$J_1,\ldots,J_n$ and real coefficients $a_1,\ldots,a_n$, with
\[
 J_i\text{ order-connected},\quad J_i\cap J_j=\varnothing\ (i\ne j),
 \qquad C(x)=\sum_{i=1}^n a_i\mathbf 1_{J_i}(x)\quad(x\in A).
\]
Order-connected subsets of $\mathbb R$ are intervals, with any
endpoint convention and possibly empty or unbounded intervals.
Lean uses the actual set indicator and an actual finite sum.

For each $x$ its membership pattern in the $J_i$ is an element of
the finite type $\{0,1\}^n$. Its function value is the sum of the
coefficients selected by that pattern. Thus the actual restricted
image $C(A)$ is a subset of the finite range of the pattern-value
function. The theorem \texttt{finite\_steps\_finite\_image} proves
this finite-image bridge. It does not postulate the image bound.

\newpage
\section{The countable interval representation}
The more permissive predicate \texttt{CountableIntervalStepOn A C}
requires actual pairwise disjoint order-connected interval sets
$J_i$ and coefficients $a_i$, indexed by $i\in\mathbb N$, such that
\[
 A\subseteq\bigcup_{i\in\mathbb N}J_i,
 \qquad C(x)=a_i\quad(x\in A\cap J_i).
\]
This includes countable interval partitions into constant pieces;
no local-finiteness assumption is needed for the obstruction.
For each $x\in A$, the actual covering condition supplies an index
$i$ with $C(x)=a_i$. Consequently $C(A)$ is contained in the actual
countable range of $i\mapsto a_i$. Lean proves this conclusion as
\texttt{countable\_steps\_countable\_image}.

\section{The actual formula has uncountable restricted image}
For $c\in I$, the strict lower bound $c>1/2$ gives $2c>1$, so
\[
 F(c)=2c\pi.
\]
If $x,y\in I$ and $F(x)=F(y)$, cancellation of the actual nonzero
real number $\pi$ gives $2x=2y$, hence $x=y$. This proves actual
injectivity on the window, using Mathlib's real $\pi$.

The real interval $I$ is uncountable. If $F(I)$ were countable,
the injective map from the actual interval subtype into its actual
image subtype would make $I$ countable. Lean proves this
contradiction explicitly. Therefore neither of the two step
representations can agree with $F$ throughout $I$.

\section{Quantifiers and the original conjunction}
The final necessary-clause theorems quantify over \emph{every}
function $C:\mathbb R\to\mathbb R$ and deny either step
representation on $I$ together with $C=F$ there. Additional theorems
formally restrict each representation from a larger set to a subset.
Since $I\subset(0,\infty)$, Lean also explicitly proves
\[
 \nexists C:\mathbb R\to\mathbb R:\quad
 C\text{ is an interval step function on }(0,\infty)
 \ \land\ \forall c>0,\ C(c)=\max(1,2c)\pi,
\]
for both the finite and countable representations. Only positive
parameters are needed; no claim about negative singular weights or
an enlarged parameter domain is used. These incompatible necessary
conditions suffice to refute the printed conjecture.

\section{Verification}
Lean 4.33.0 and Mathlib commit
\begin{center}\small\texttt{db584cd6d46c92f209a44c0f1c829460d327499d}\end{center}
are pinned. The package includes a new root build, transitive axiom
audits of all 15 theorems and full official
\texttt{leanchecker --fresh Main} replay. No \texttt{sorry},
\texttt{native\_decide} or additional axiom is used. OpenAI Codex
assisted; a separate supervising Codex agent reviewed source,
report and PDF. This is AI review, not human certification.
Organizer acceptance is pending.
\begin{thebibliography}{1}
\bibitem{Ryzhik} L.~Ryzhik, Stanford Math 205A, Fall 2015,
Homework 3 solutions, Problem 1(i), p.~1.
\url{https://math.stanford.edu/~ryzhik/STANFORD/STANF205-16/205A_hw3_sol.pdf}.
\end{thebibliography}
\end{document}
